\documentclass[12pt,a4paper]{article}
\usepackage[utf8]{inputenc}
\usepackage[T5]{fontenc}
\usepackage[vietnamese]{babel}
\usepackage[margin=2.5cm]{geometry}

\title{Mô tả thư mục checkpoint}
\date{}

\begin{document}
\maketitle

Thư mục \texttt{checkpoints} dùng để lưu các tệp trạng thái mô hình đã được lưu trong quá trình huấn luyện. Với bài tập MLP, các tệp này thường có phần mở rộng \texttt{.pth} hoặc \texttt{.pt} và chứa trọng số, tham số đã học của mạng nơ-ron; tùy cách cài đặt, chúng cũng có thể lưu thêm thông tin như epoch tốt nhất hoặc trạng thái bộ tối ưu.

Nhờ checkpoint, có thể nạp lại mô hình đã huấn luyện mà không cần huấn luyện lại từ đầu. Đây là các tệp đầu ra được tạo khi chạy chương trình, không phải mã nguồn Python.

\end{document}
